% !PRINTED 646
\weberpage{646}
Bedeutet $t$ eine neue Variable, so ist
\[
N(\delta t - \delta_0) = D E(x t - x_0,\ y t - y_0, \ldots),
\]
also, wenn man nach absteigenden Potenzen von $t$ ordnet und den Faktor $N(\delta) = D E$ beiderseits forthebt:
\weq{4}{N\left(t - \frac{\delta_0}{\delta}\right) = t^n + \overline{C}_1 t^{n-1} + \overline{C}_2 t^{n-2} + \cdots,}
worin $\overline{C}_1, \overline{C}_2 \ldots$ keinen anderen Nenner als $E$ haben, und folglich ganze Funktionale in $\overline{Z}$ sind. Demnach ist nach der Definition \S\,179, 6. auch $\delta_0 : \delta$ ein ganzes Funktional, wie bewiesen werden sollte.

Demnach erhalten wir folgende Definitionen:

1. \gesperrt{Das Funktional $\delta = \alpha x + \beta y + \cdots$ und jedes mit $\delta$ assoziierte Funktional ist der größte gemeinschaftliche Teiler von $\alpha, \beta, \ldots$}

2. \gesperrt{Ist $\alpha x + \beta y$ eine Einheit, so heißen $\alpha$ und $\beta$ teilerfremd oder relativ prim. Gibt es Funktionale $\xi, \eta$, für die $\alpha \xi + \beta \eta$ eine Einheit ist, so sind $\alpha, \beta$ relativ prim.}

3. \gesperrt{Ist $\alpha$ relativ prim zu $\beta$ und zu $\gamma$, so ist es auch relativ prim zu $\beta\gamma$.}

Denn nach Voraussetzung sind
\[
\varepsilon = \alpha x + \beta y,\quad \varepsilon_1 = \alpha u + \beta v
% [sic] printed epsilon_1 = alpha u + beta v, should be alpha u + gamma v: alpha is assumed prime to gamma, and the next display expands (alpha x + beta y)(alpha u + gamma v)
\]
Einheiten ($x, y, u, v$ neue Variablen).
\[
\alpha(\alpha u x + \gamma v x + \beta u y) + \beta \gamma v y = \varepsilon \varepsilon_1,
\]
also nach 3. $\alpha$ und $\beta\gamma$ relativ prim.
% [sic] printed "nach 3.", a circular reference to statement 3. itself, which is being proved; the step uses the criterion of 2. (alpha xi + beta eta a unit)

4. \gesperrt{Ist $\alpha$ relativ prim zu $\beta$ und $\alpha\mu$ durch $\beta$ teilbar, so ist $\mu$ durch $\beta$ teilbar.}

Denn aus der Voraussetzung folgt:
\[
\alpha \mu x + \beta \mu y = \varepsilon \mu,
\]
woraus, da $\varepsilon$ eine Einheit ist, sich der Beweis ergibt.

\begin{center}
\S\,182. \textbf{Primfunktionale in $\overline{\Omega}$.}
\end{center}

1. \gesperrt{Definition. Ein ganzes Funktional $\pi$ in $\overline{\Omega}$ heißt ein Primfunktional, wenn es keine Einheit ist und außer durch die Einheiten nur durch die mit ihm assoziierten Funktionale teilbar ist.}
